\documentclass[11pt]{article}

\usepackage[a4paper,margin=2.2cm]{geometry}
\usepackage[T1]{fontenc}
\usepackage[utf8]{inputenc}
\usepackage{lmodern}
\usepackage{amsmath,amssymb,mathtools,bm}
\usepackage{physics}
\usepackage{microtype}
\usepackage[dvipsnames]{xcolor}
\usepackage[most]{tcolorbox}
\usepackage{enumitem}
\usepackage{titlesec}
\usepackage{booktabs}

% ---------- Section formatting ----------
\titleformat{\section}{\Large\bfseries}{\thesection}{0.75em}{}
\titleformat{\subsection}{\large\bfseries}{\thesubsection}{0.75em}{}
\titlespacing*{\section}{0pt}{1.4ex plus .2ex minus .2ex}{0.9ex}
\titlespacing*{\subsection}{0pt}{1.0ex plus .2ex minus .2ex}{0.5ex}

% ---------- Box environments ----------
\newtcolorbox{resultbox}[1]{
    breakable, enhanced,
    colback=blue!2, colframe=blue!40!black,
    boxrule=0.7pt, arc=1.5mm,
    left=3mm, right=3mm, top=1.5mm, bottom=1.5mm,
    title={\bfseries #1}, coltitle=black,
    colbacktitle=blue!10, fonttitle=\bfseries
}

\newtcolorbox{notebox}[1]{
    breakable, enhanced,
    colback=gray!4, colframe=gray!50,
    boxrule=0.7pt, arc=1.5mm,
    left=3mm, right=3mm, top=1.5mm, bottom=1.5mm,
    title={\bfseries #1}, coltitle=black,
    colbacktitle=gray!12, fonttitle=\bfseries
}

% ---------- Source note ----------
\newcommand{\sourcenote}[1]{%
    \par\vspace{0.4em}%
    \textcolor{gray}{\footnotesize\textbf{Source:} #1}%
}

% ---------- Math shortcuts ----------
\newcommand{\R}{\mathbb{R}}

\begin{document}

\begin{center}
    {\LARGE\bfseries LPV-LFR Realization of the Gantry First-Principles Model}\\[0.5em]
    {\large Derivation from the MATLAB CT State-Space}\\[0.4em]
    {\normalsize Baseline for the LFR Augmentation Framework}
\end{center}

\vspace{1em}

%------------------------------------------------------------------------------
\section{Starting Point: CT State-Space from MATLAB}
%------------------------------------------------------------------------------

The gantry first-principles model is derived from the Euler-Lagrange equations
for the dual-gantry stage in García-Herreros et al. \cite{garcia2013model}.
In logical coordinates $q = [X,\,\Theta,\,Y]^\top \in \R^3$, the equations of motion are
\begin{equation}
    M(Y)\,\ddot{q}(t) + C\,\dot{q}(t) + K\,q(t) = f_\ell(t),
    \label{eq:eom}
\end{equation}
with
\begin{equation}
    M(Y) =
    \begin{bmatrix}
        m_1 + m_2 + m_b + m_h &
        \dfrac{(m_1 - m_2)L_b}{2} - m_h Y &
        0 \\[6pt]
        \dfrac{(m_1 - m_2)L_b}{2} - m_h Y &
        J_b + J_h + \dfrac{(m_1+m_2)L_b^2}{4} + m_h d^2 + m_h Y^2 &
        -m_h d \\[6pt]
        0 & -m_h d & m_h
    \end{bmatrix},
    \label{eq:M_full}
\end{equation}
\begin{equation}
    C =
    \begin{bmatrix}
        c_{g1} + c_{g2} & \dfrac{(c_{g1} - c_{g2})L_b}{2} & 0 \\[6pt]
        \dfrac{(c_{g1} - c_{g2})L_b}{2} &
        c_{b1} + c_{b2} + \dfrac{(c_{g1}+c_{g2})L_b^2}{4} &
        0 \\[6pt]
        0 & 0 & c_y
    \end{bmatrix},
    \qquad
    K =
    \begin{bmatrix}
        0 & 0 & 0 \\
        0 & k_{b1} + k_{b2} & 0 \\
        0 & 0 & 0
    \end{bmatrix},
    \label{eq:CK_full}
\end{equation}
and
\begin{equation}
    f_\ell =
    \begin{bmatrix}
        F_X \\
        F_\Theta \\
        F_Y
    \end{bmatrix}
    \in \R^3
    \label{eq:fell}
\end{equation}
the generalized-force vector in logical coordinates.

If the actuator forces are instead given in stage coordinates as
\begin{equation}
    f_s =
    \begin{bmatrix}
        F_{X1} \\
        F_{X2} \\
        F_Y
    \end{bmatrix},
    \label{eq:fs}
\end{equation}
then
\begin{equation}
    f_\ell = P\,f_s,
    \qquad
    P =
    \begin{bmatrix}
        1 & 1 & 0 \\
        L_b/2 & -L_b/2 & 0 \\
        0 & 0 & 1
    \end{bmatrix},
    \label{eq:P_force}
\end{equation}
which is the constant input transformation used in the MATLAB/Simulink model.

From this point onward, we denote the generalized input by
\begin{equation}
    u := f_\ell
    \label{eq:u_def}
\end{equation}
to match standard state-space and LFR notation.

Introducing the state $x = [q^\top,\,\dot{q}^\top]^\top \in \R^6$, the CT state-space is
\begin{equation}
    \dot{x}(t) = A_c(Y)\,x(t) + B_c(Y)\,u(t),
    \qquad
    y(t) = C_c\,x(t),
    \label{eq:ctss}
\end{equation}
with
\begin{equation}
    A_c(Y) = \begin{bmatrix} 0 & I_3 \\ -M(Y)^{-1}K & -M(Y)^{-1}C \end{bmatrix},
    \quad
    B_c(Y) = \begin{bmatrix} 0 \\ M(Y)^{-1} \end{bmatrix},
    \quad
    C_c = \begin{bmatrix} I_3 & 0 \end{bmatrix}.
    \label{eq:ctmatrices}
\end{equation}
The parameter dependence is rational in $Y$ through $M(Y)^{-1}$.
Consequently, the model is not in a standard affine LPV state-space form in the
scheduling variable $Y$.
The construction below gives an exact LPV-LFR realization of this rational dependence.

\sourcenote{\texttt{kamtin-fp-model/03 Simulink gantry/functions/getss.m}}

\begin{notebox}{Stage coordinates and implementation}
The derivation below is written in logical coordinates
$q = [X,\,\Theta,\,Y]^\top$, because this is the coordinate system in which the
MATLAB equations are defined.

If the implemented model must operate in stage coordinates
$q_s = [q_{X1},\,q_{X2},\,q_Y]^\top$, then $q_s = P^\top q$ and the full state
must be transformed by the corresponding block transformation
$x_s = \mathrm{blkdiag}(P^\top,\,P^\top)\,x$.
In that case, the realized baseline can be similarity-transformed to stage
coordinates before implementation.

If only the measured output is needed in stage coordinates, then the output row
becomes $C_y = P^\top [I_3,\,0]$ instead of $[I_3,\,0]$.
\end{notebox}

%------------------------------------------------------------------------------
\section{LPV-LFR Definition}
%------------------------------------------------------------------------------

An LPV system admits a Linear Fractional Representation (LFR)
when it can be written as the lower linear fractional transformation
$\mathcal{F}_\ell(G,\,\Delta(p))$,
where $G$ is a constant (parameter-independent) matrix partitioned as
\begin{equation}
    G = \begin{bmatrix}
        A   & B_w   & B_u \\
        C_z & D_{zw} & D_{zu} \\
        C_y & D_{yw} & D_{yu}
    \end{bmatrix},
    \label{eq:Gmat}
\end{equation}
and $\Delta(p)$ is a block-diagonal matrix whose entries depend only on the
scheduling parameter $p$.

The system equations are
\begin{align}
    \dot{x} &= A\,x + B_w\,w + B_u\,u, \label{eq:lfr_state}\\
    z       &= C_z\,x + D_{zw}\,w + D_{zu}\,u, \label{eq:lfr_z}\\
    y       &= C_y\,x + D_{yw}\,w + D_{yu}\,u, \label{eq:lfr_out}\\
    w       &= \Delta(p)\,z. \label{eq:lfr_loop}
\end{align}
Substituting \eqref{eq:lfr_loop} into \eqref{eq:lfr_state}--\eqref{eq:lfr_out} recovers the
original parameter-dependent state-space, provided the algebraic loop
\eqref{eq:lfr_z}--\eqref{eq:lfr_loop} has a unique solution for $z$ given $x$ and $u$
(well-posedness, see Section~\ref{sec:wp}).

\sourcenote{Primary CT source: Drenth thesis \cite{drenth2025thesis},
Section~2.1, eq.\ (2.1)--(2.3). Supporting DT companion: Drenth et al.,
IFAC 2025 \cite{drenth2025lpvlfr}, eq.\ (6)--(9).}

%------------------------------------------------------------------------------
\section{Decomposition of \texorpdfstring{$M(Y)$}{M(Y)}}
%------------------------------------------------------------------------------

The inertia matrix in \eqref{eq:M_full} is polynomial in $Y$ of degree~2.
Rearranging the already stated expression into terms of equal degree in $Y$ gives:
\begin{equation}
    M(Y) = \underbrace{\begin{bmatrix} a & b_0 & 0 \\ b_0 & c_0 & -e \\ 0 & -e & f \end{bmatrix}}_{M_0}
         + Y \underbrace{\begin{bmatrix} 0 & -m_h & 0 \\ -m_h & 0 & 0 \\ 0 & 0 & 0 \end{bmatrix}}_{M_1}
         + Y^2 \underbrace{\begin{bmatrix} 0 & 0 & 0 \\ 0 & m_h & 0 \\ 0 & 0 & 0 \end{bmatrix}}_{M_2},
    \label{eq:decomp}
\end{equation}
where $b_0 = \tfrac{(m_1-m_2)L_b}{2}$ and
$c_0 = J_b + J_h + \tfrac{(m_1+m_2)L_b^2}{4} + m_h d^2$.
Compactly:
\begin{equation}
    M(Y) = M_0 + Y\,M_1 + Y^2\,M_2, \qquad M_0 = M(0).
    \label{eq:decomp_short}
\end{equation}
$M_1$ is rank~2 (only the $(1,2)$ and $(2,1)$ entries are nonzero);
$M_2$ is rank~1 (only the $(2,2)$ entry is nonzero).
The choice $M_0 = M(0)$ is algebraically convenient; numerically any fixed
$M_0 = M(Y_0)$ is equivalent.

%------------------------------------------------------------------------------
\section{Latent Variable Construction}
%------------------------------------------------------------------------------

Let $f_\mathrm{gen} = [-K,\,-C]\,x + u$ denote the generalized force.
Define the latent acceleration vector
\begin{equation}
    v \;=\; M(Y)^{-1}\,f_\mathrm{gen} \;=\; \ddot{q} \;\in\; \R^3.
    \label{eq:latent_v}
\end{equation}
The $Y$-dependence in $M(Y)^{-1}$ enters through two additional latent variables:
\begin{equation}
    v_1 \;=\; Y\,v, \qquad v_2 \;=\; Y^2\,v \;=\; Y\,v_1.
    \label{eq:latent_v1v2}
\end{equation}
Stack these into loop signals
\begin{equation}
    z \;=\; \begin{bmatrix} v \\ v_1 \end{bmatrix} \in \R^6,
    \qquad
    w \;=\; \begin{bmatrix} v_1 \\ v_2 \end{bmatrix} \in \R^6.
    \label{eq:zw}
\end{equation}
Observe that $w = Y\,z$, so define the scheduling block
\begin{equation}
    \Delta(Y) \;=\; Y\,I_6, \qquad n_z = n_w = 6.
    \label{eq:Delta}
\end{equation}
The scalar $Y$ enters the model exclusively through $\Delta(Y)$;
all remaining structure is parameter-independent.
This establishes a valid LPV-LFR realization of the baseline model.
No claim of minimality is made for the chosen latent dimension.

%------------------------------------------------------------------------------
\section{Constant \texorpdfstring{$G$}{G} Matrix}
%------------------------------------------------------------------------------

The submatrices of $G$ are determined by matching \eqref{eq:lfr_state}--\eqref{eq:lfr_out}
to the CT state-space \eqref{eq:ctss}--\eqref{eq:ctmatrices}.

\subsection{State equation}

The state equation $\dot{x} = A_c(Y)\,x + B_c(Y)\,u$ has two blocks:
\begin{align}
    \dot{q}   &= \dot{q}, \label{eq:upper}\\
    \ddot{q}  &= -M(Y)^{-1}K\,q - M(Y)^{-1}C\,\dot{q} + M(Y)^{-1}u = v. \label{eq:lower}
\end{align}
The upper block depends only on $x$ and contributes $[0,\,I_3;\,0,\,0]\,x$.
The lower block is $v = M(Y)^{-1}f_\mathrm{gen}$.

To express $v$ in terms of $x$, $w$, and $u$, pre-multiply definition \eqref{eq:latent_v}
by $M(Y)$ and expand using \eqref{eq:decomp_short}:
\begin{equation}
    M_0\,v \;=\; f_\mathrm{gen} - Y\,M_1\,v - Y^2\,M_2\,v
              \;=\; f_\mathrm{gen} - M_1\,v_1 - M_2\,v_2
              \;=\; [-K,-C]\,x - M_1\,w_1 - M_2\,w_2 + u,
    \label{eq:Mv_expand}
\end{equation}
where $w_1 = v_1$ and $w_2 = v_2$ denote the $\R^3$ blocks of $w$.
Solving for $v$:
\begin{equation}
    v \;=\; M_0^{-1}[-K,-C]\,x - M_0^{-1}M_1\,w_1 - M_0^{-1}M_2\,w_2 + M_0^{-1}u.
    \label{eq:v_expr}
\end{equation}
Substituting into the lower block of the state equation:
\begin{equation}
    \ddot{q} \;=\;
    \bigl(-M_0^{-1}K\,q - M_0^{-1}C\,\dot{q}\bigr)
    - M_0^{-1}M_1\,w_1 - M_0^{-1}M_2\,w_2
    + M_0^{-1}u.
    \label{eq:lower_expanded}
\end{equation}
Reading off $A$, $B_w$, $B_u$ from \eqref{eq:upper} and \eqref{eq:lower_expanded}:
\begin{equation}
    A = \begin{bmatrix} 0 & I_3 \\ -M_0^{-1}K & -M_0^{-1}C \end{bmatrix},
    \quad
    B_w = \begin{bmatrix} 0 & 0 \\ -M_0^{-1}M_1 & -M_0^{-1}M_2 \end{bmatrix},
    \quad
    B_u = \begin{bmatrix} 0 \\ M_0^{-1} \end{bmatrix}.
    \label{eq:state_mats}
\end{equation}
For the choice $M_0 = M(0)$, the constant matrix $A$ equals $A_c(0)$.
More generally, if a different constant reference matrix $M_0 = M(Y_0)$ were chosen,
the same construction would proceed with that choice.

\subsection{Loop equation}

The loop equation \eqref{eq:lfr_z} must produce $z = [v;\,v_1]$.
From \eqref{eq:v_expr}, the first block is $v$ expressed in terms of $x$, $w$, $u$.
The second block is $v_1 = w_1 = I_3\,w_1 + 0\cdot w_2$.
Stacking both blocks:
\begin{equation}
    z \;=\;
    \underbrace{\begin{bmatrix} M_0^{-1}[-K,\,-C] \\ 0 \end{bmatrix}}_{C_z}\,x
    \;+\;
    \underbrace{\begin{bmatrix} -M_0^{-1}M_1 & -M_0^{-1}M_2 \\ I_3 & 0 \end{bmatrix}}_{D_{zw}}\,w
    \;+\;
    \underbrace{\begin{bmatrix} M_0^{-1} \\ 0 \end{bmatrix}}_{D_{zu}}\,u.
    \label{eq:z_eq}
\end{equation}

\subsection{Output equation}

The output in logical coordinates is $y = C_c\,x = [I_3,\,0]\,x$ with no feedthrough:
\begin{equation}
    C_y = \begin{bmatrix} I_3 & 0 \end{bmatrix}, \qquad D_{yw} = 0, \qquad D_{yu} = 0.
    \label{eq:out_eq}
\end{equation}

\medskip
\begin{resultbox}{Complete $G$ matrix}
With $z = [v;\,v_1]$, $w = [v_1;\,v_2]$, $\Delta(Y) = Y\,I_6$:
\begin{alignat}{2}
    A      &= \begin{bmatrix} 0 & I_3 \\ -M_0^{-1}K & -M_0^{-1}C \end{bmatrix}
             & \;\in\; \R^{6\times6}, \label{eq:A_fin}\\[4pt]
    B_w    &= \begin{bmatrix} 0 & 0 \\ -M_0^{-1}M_1 & -M_0^{-1}M_2 \end{bmatrix}
             & \;\in\; \R^{6\times6}, \label{eq:Bw_fin}\\[4pt]
    B_u    &= \begin{bmatrix} 0 \\ M_0^{-1} \end{bmatrix}
             & \;\in\; \R^{6\times3}, \label{eq:Bu_fin}\\[4pt]
    C_z    &= \begin{bmatrix} M_0^{-1}[-K,\,-C] \\ 0 \end{bmatrix}
             & \;\in\; \R^{6\times6}, \label{eq:Cz_fin}\\[4pt]
    D_{zw} &= \begin{bmatrix} -M_0^{-1}M_1 & -M_0^{-1}M_2 \\ I_3 & 0 \end{bmatrix}
             & \;\in\; \R^{6\times6}, \label{eq:Dzw_fin}\\[4pt]
    D_{zu} &= \begin{bmatrix} M_0^{-1} \\ 0 \end{bmatrix}
             & \;\in\; \R^{6\times3}, \label{eq:Dzu_fin}\\[4pt]
    C_y    &= \begin{bmatrix} I_3 & 0 \end{bmatrix}, \quad D_{yw} = 0, \quad D_{yu} = 0.
             & \label{eq:out_fin}
\end{alignat}
All entries depend only on $M_0, M_1, M_2, K, C$ (all constant matrices).
\end{resultbox}

\begin{notebox}{Representation versus runtime evaluation}
The matrices above define a constant LPV-LFR representation of the baseline.
For simulation or training, one may either:
(i) evaluate the explicit LFR algebraic loop at each scheduling value $Y$, or
(ii) use the equivalent collapsed state equation
$\dot{x} = A_c(Y)x + B_c(Y)u$ obtained after solving the loop.
These two descriptions are algebraically equivalent for the present construction.
\end{notebox}

%------------------------------------------------------------------------------
\section{Algebraic Verification}
\label{sec:verify}
%------------------------------------------------------------------------------

We verify that substituting $w = \Delta(Y)\,z = Y\,z$ into the LFR equations
recovers \eqref{eq:ctss}--\eqref{eq:ctmatrices}.

\medskip
\noindent\textbf{Step 1: Solve the loop.}
With $w = Y\,z$, i.e., $w_1 = Y\,v$ and $w_2 = Y\,v_1 = Y^2\,v$,
the first block of \eqref{eq:z_eq} becomes:
\begin{align}
    v &= M_0^{-1}[-K,-C]\,x
        - M_0^{-1}M_1\,(Y\,v)
        - M_0^{-1}M_2\,(Y^2\,v)
        + M_0^{-1}u \notag\\
    M_0\,v + Y\,M_1\,v + Y^2\,M_2\,v &= [-K,-C]\,x + u \notag\\
    M(Y)\,v &= f_\mathrm{gen}. \label{eq:loop_check}
\end{align}
This is exactly \eqref{eq:latent_v}. The second block $z[3{:}6] = v_1 = w_1 = Y\,v$
is trivially consistent with $v_1 = Y\,v$.

\medskip
\noindent\textbf{Step 2: Recover the state equation.}
Substituting $w = [Y\,v;\,Y^2\,v]$ into \eqref{eq:lfr_state} with \eqref{eq:state_mats}:
\begin{align}
    \dot{x} &= \begin{bmatrix} 0 & I_3 \\ -M_0^{-1}K & -M_0^{-1}C \end{bmatrix} x
              + \begin{bmatrix} 0 & 0 \\ -M_0^{-1}M_1 & -M_0^{-1}M_2 \end{bmatrix}
                \begin{bmatrix} Y\,v \\ Y^2\,v \end{bmatrix}
              + \begin{bmatrix} 0 \\ M_0^{-1} \end{bmatrix} u \notag\\[4pt]
    &= \begin{bmatrix} \dot{q} \\
         -M_0^{-1}K\,q - M_0^{-1}C\,\dot{q}
         - M_0^{-1}M_1\,Y\,v - M_0^{-1}M_2\,Y^2\,v
         + M_0^{-1}u \end{bmatrix} \notag\\[4pt]
    &= \begin{bmatrix} \dot{q} \\
         M_0^{-1}\bigl(f_\mathrm{gen} - M_1\,Y\,v - M_2\,Y^2\,v\bigr) \end{bmatrix}.
    \label{eq:xdot_expand}
\end{align}
From \eqref{eq:loop_check}: $M_0\,v = f_\mathrm{gen} - M_1\,Y\,v - M_2\,Y^2\,v$.
Substituting:
\begin{equation}
    \dot{x} = \begin{bmatrix} \dot{q} \\ M_0^{-1}\,M_0\,v \end{bmatrix}
            = \begin{bmatrix} \dot{q} \\ v \end{bmatrix}
            = \begin{bmatrix} \dot{q} \\ M(Y)^{-1}f_\mathrm{gen} \end{bmatrix}
            = A_c(Y)\,x + B_c(Y)\,u. \quad \checkmark
    \label{eq:verify_ok}
\end{equation}

\medskip
\noindent\textbf{Step 3: Output.}
$y = C_y\,x = [I_3,\,0]\,x = q = C_c\,x$. \quad $\checkmark$

%------------------------------------------------------------------------------
\section{Well-Posedness of the Algebraic Loop}
\label{sec:wp}
%------------------------------------------------------------------------------

The LFR interconnection is well-posed if the algebraic loop
\eqref{eq:lfr_z}--\eqref{eq:lfr_loop} has a unique solution for $z$ given any
$(x, u)$ and any scheduling value $Y$.
The general condition is
\begin{equation}
    \det\bigl(I - D_{zw}\,\Delta(Y)\bigr) \neq 0 \qquad \forall\,Y.
    \label{eq:wp_general}
\end{equation}

\sourcenote{Drenth thesis \cite{drenth2025thesis}, Section~2.2, eq.\ (2.4).}

\medskip
\noindent\textbf{Reduction to $M(Y)$ invertibility.}
With $\Delta(Y) = Y\,I_6$, the loop equation \eqref{eq:z_eq} with $w = Y\,z$ is:
\begin{align}
    v   &= M_0^{-1}[-K,-C]\,x - M_0^{-1}M_1\,(Y\,v) - M_0^{-1}M_2\,(Y\,v_1) + M_0^{-1}u,
    \label{eq:loop_v}\\
    v_1 &= Y\,v. \label{eq:loop_v1}
\end{align}
Substituting \eqref{eq:loop_v1} into \eqref{eq:loop_v}:
\begin{equation}
    \bigl(M_0 + Y\,M_1 + Y^2\,M_2\bigr)\,v = f_\mathrm{gen}
    \quad\Longrightarrow\quad
    M(Y)\,v = f_\mathrm{gen}.
    \label{eq:loop_reduced}
\end{equation}
The $6\times6$ algebraic loop reduces to a $3\times3$ linear system for $v$.
Once $v$ is found, $v_1 = Y\,v$ follows immediately.
The loop has a unique solution if and only if $M(Y)$ is invertible.

\medskip
\begin{notebox}{Why this reduction is specific to this $D_{zw}$}
    The collapse \eqref{eq:loop_reduced} is not a general property of LFRs.
    It holds because $D_{zw}$ was constructed so that pre-multiplying by $M_0$ and collecting
    terms in $v$ reconstitutes exactly $M(Y)$ on the left-hand side.
    For a different $D_{zw}$, the loop would not reduce in this way.
    The implication ``loop has unique solution $\Leftrightarrow$ $M(Y)$ invertible'' is
    a consequence of this specific construction, not a general theorem.
\end{notebox}

\medskip
\noindent\textbf{Invertibility of $M(Y)$.}
The companion document `LPV/supporting/derivations/M-invertibility.tex` proves via Sylvester's criterion that
$M(Y)$ is positive definite for all $Y \in \R$, provided
$m_1, m_2, m_b, m_h > 0$, $J_b, J_h > 0$, $L_b, d > 0$.
Positive definiteness implies invertibility.

\medskip
\begin{resultbox}{Well-posedness}
    The algebraic loop of the LFR has a unique solution for every $Y \in \R$ and every
    $(x, u)$, provided all physical mass, inertia, and geometry parameters are
    strictly positive.
    Well-posedness holds globally (not only on the operational range $Y \in [0.05,\,0.75]$~m).
\end{resultbox}

\newpage
\begin{thebibliography}{9}

\bibitem{garcia2013model}
I.~Garc{\'i}a-Herreros, X.~Kestelyn, J.~Gomand, R.~Coleman, and P.-J.~Barre,
``Model-based decoupling control method for dual-drive gantry stages:
A case study with experimental validations,''
\textit{Control Engineering Practice}, vol.~21, no.~3, pp.~298--307, 2013.
doi: 10.1016/j.conengprac.2012.10.010.

\bibitem{drenth2025lpvlfr}
J.~Drenth, J.~Hoekstra, M.~Schoukens, and R.~T{\'o}th,
``Efficient Learning of Affine and Rational Dependency LPV Models With LFR,''
in \textit{Proc.\ IFAC}, 2025.

\bibitem{drenth2025thesis}
J.~Drenth,
``Gradient-based Learning of LPV-LFR Models,''
M.Sc.\ thesis, Eindhoven University of Technology, 2025.

\end{thebibliography}

\end{document}
